### Conventions

* All graphs are **finite, simple, undirected**.
* For graphs (A,B), the (two–color) **Ramsey number** (R(A,B)) is the least integer (N) such that for every red/blue coloring of the edges of the complete graph (K_N), either the red graph contains a (not-necessarily-induced) copy of (A) or the blue graph contains a copy of (B). Equivalently: for every graph (X) on (N) vertices, either (X) contains (A) or (\overline X) contains (B).
* “(\ll m)” means “(\le C,m)” for some constant (C) depending only on (G), not on (H) or (m).

### The fixed graphs (G)

Let (G) be one of:

* (Q_3): the 3-dimensional cube graph (8 vertices, 12 edges).
* (K_{3,3}): complete bipartite graph with parts of size 3 (6 vertices, 9 edges).
* (H_5): the graph obtained from (C_5) by adding two vertex-disjoint chords; equivalently (K_4^*), the graph obtained from (K_4) by subdividing one edge once.

### The variable graph (H)

Let (H) be any graph with

* (e(H)=m) edges,
* **no isolated vertices** (i.e., (\delta(H)\ge 1)).

Note: the “no isolated vertices” condition implies (v(H)\le 2m).

### Question interpreted as a statement

For each (G\in{Q_3,K_{3,3},H_5}), is it true that

[
\exists C_G>0\ \forall H\ (e(H)=m,\ \delta(H)\ge 1)\ \Longrightarrow\ R(G,H)\le C_G,m\ ?
]

If true, we say (G) is **Ramsey size-linear**.

---
